% Part: first-order-logic
% Chapter: model-theory
% Section: isomorphism

\documentclass[../../../include/open-logic-section]{subfiles}

\begin{document}

\olfileid{mod}{bas}{iso}
\olsection{समरूपी रचना}

प्रथम-क्रमाच्या !!{structure}s दोन प्रकारे एकसारख्या असू शकतात. त्यांच्या
सारखेपणाचा एक प्रकार असा की त्या तीच !!{sentence}s सत्य ठरवतात. अशा
!!{structure}s ना आपण \emph{प्राथमिक सममूल्य} म्हणतो. परंतु रचना अतिशय
भिन्न असूनही तीच !!{sentence}s सत्य ठरवू शकतात---उदाहरणार्थ, एक रचना
!!{enumerable} असू शकते आणि दुसरी नसू शकते. याचे कारण असे की
!!a{structure} ची पुष्कळ वैशिष्ट्ये प्रथम-क्रम भाषांमध्ये व्यक्त करता येत
नाहीत: कधी भाषा पुरेशी समृद्ध नसते, तर कधी लोव्हेनहाइम--स्कोलेम
प्रमेयासारख्या प्रथम-क्रम तर्कशास्त्राच्या मूलभूत मर्यादा आड येतात. म्हणून
!!{structure}s च्या सारखेपणाचा दुसरा, अधिक काटेकोर प्रकार असा की त्या
मूलतः एकच असतात: त्यांना घडवणाऱ्या वस्तू वेगळ्या असतात, पण त्यांची
रचनात्मक वैशिष्ट्ये वेगळी नसतात. \emph{समरूपते} च्या संकल्पनेने हे नेमके
मांडता येते.

\begin{defn}
\ollabel{defn:elem-equiv}
एकाच !!{language}~$\Lang{L}$ साठीच्या दोन !!{structure}s $\Struct{M}$ आणि
$\Struct M'$ दिलेल्या असताना, $\Lang{L}$ मधील प्रत्येक !!{sentence}~$!A$
साठी $\Sat{M}{!A}$ तेव्हा आणि केवळ तेव्हाच जेव्हा $\Sat{M'}{!A}$ असेल,
तर $\Struct{M}$ ही $\Struct M'$ शी \emph{प्राथमिक सममूल्य} आहे असे म्हणतो
आणि $\Struct{M} \equiv \Struct M'$ असे लिहितो.
\end{defn}

\begin{defn}
\ollabel{defn:isomorphism} एकाच !!{language}~$\Lang L$ साठीच्या दोन
!!{structure}s $\Struct{M}$ आणि $\Struct M'$ दिलेल्या असताना, पुढील अटी
पूर्ण करणारे $h \colon \Domain{M} \to \Domain{M'}$ फलन असेल, तर
$\Struct{M}$ ही~$\Struct M'$ शी \emph{समरूपी} आहे असे म्हणतो आणि
$\Struct{M} \simeq \Struct M'$ असे लिहितो:
\begin{enumerate}
\item $h$ हे !!{injective} आहे: $h(x) =
  h(y)$ असेल, तर $x = y$;
\item $h$ हे !!{surjective} आहे: प्रत्येक $y \in \Domain{M'}$ साठी
  $h(x) = y$ असा $x \in \Domain{M}$ असतो;
\item \ollabel{defn:iso-const}प्रत्येक !!{constant} $c$ साठी:
  $h(\Assign{c}{M}) = \Assign{c}{M'}$;
\item \ollabel{defn:iso-pred}प्रत्येक $n$-स्थानी !!{predicate}~$P$ साठी:
  \[
  \tuple{a_1, \dots, a_n}\in \Assign{P}{M} \quad\text{तेव्हा आणि केवळ तेव्हाच जेव्हा}\quad
  \tuple{h(a_1), \dots, h(a_n)} \in \Assign{P}{M'};
  \]
\item \ollabel{defn:iso-func}प्रत्येक $n$-स्थानी !!{function} $f$ साठी:
  \[
  h(\Assign{f}{M}(a_1, \dots, a_n)) =
  \Assign{f}{M'}(h(a_1), \dots, h(a_n)).
  \]
\end{enumerate}
\end{defn}

\begin{thm}
\ollabel{thm:isom}
$\Struct{M} \iso \Struct M'$ असेल, तर $\Struct{M} \elemequiv
\Struct{M'}$.
\end{thm}

\begin{proof}
$h$ हे $\Struct{M}$ पासून $\Struct M'$ वरचे समरूपण असू द्या. कोणत्याही
मूल्यांकन~$s$ साठी $h \circ s$ हे $h$ आणि $s$ यांचे संयोजन असते, म्हणजे
ते $\Struct{M'}$ मधील असे मूल्यांकन आहे की $(h \circ s)(x) = h(s(x))$.
$t$ आणि $!A$ यांवरील विगमनाने पुढील अधिक सबळ दावे सिद्ध करता येतात:
\begin{enumerate}
  \item[a.] $h(\Value{t}{M}[s]) = \Value{t}{M'}[h\circ s]$.
  \item[b.] $\Sat{M}{!A}[s]$ तेव्हा आणि केवळ तेव्हाच जेव्हा $\Sat{M'}{!A}[h \circ s]$.
\end{enumerate}
पहिला दावा~$t$ च्या गुंतागुंतीवरील विगमनाने सिद्ध होतो.
\begin{enumerate}
\item $t \ident c$ असेल, तर $\Value{c}{M}[s] = \Assign{c}{M}$ आणि
  $\Value{c}{M'}[h \circ s] = \Assign{c}{M'}$. म्हणून
  $h(\Value{t}{M}[s]) = h(\Assign{c}{M}) = \Assign{c}{M'}$
  (\olref{defn:isomorphism} मधील \olref{defn:iso-const} नुसार) $=
  \Value{t}{M'}[h \circ s]$.
\item $t \ident x$ असेल, तर $\Value{x}{M}[s] = s(x)$ आणि
  $\Value{x}{M'}[h \circ s] = h(s(x))$. म्हणून $h(\Value{x}{M}[s]) =
  h(s(x)) = \Value{x}{M'}[h \circ s]$.
\item $t \ident f(t_1, \dots, t_n)$ असेल, तर
  \begin{align*}
    \Value{t}{M}[s] & = \Assign{f}{M}(\Value{t_1}{M}[s], \dots, \Value{t_n}{M}[s]) \quad\text{आणि}\\
  \Value{t}{M'}[h \circ s] & = \Assign{f}{M'}(\Value{t_1}{M'}[h \circ
    s], \dots, \Value{t_n}{M'}[h \circ s]).
  \end{align*}
  %READERNOTE{OLFOL-026 — गोठवलेल्या इंग्रजी सूत्रात M-प्राइममधील पदमूल्य काढताना f चे अर्थनिर्धारण चुकून M मधील दिले आहे. व्याख्या आणि पुढील गणना यांनुसार येथे M-प्राइम असणे आवश्यक आहे; मराठी सूत्रात ती खूण दुरुस्त केली आहे.}
  प्रत्येक $i$ साठी $h(\Value{t_i}{M}[s])
  = \Value{t_i}{M'}[h\circ s]$ हे विगमन गृहीतक आहे. म्हणून
  \begin{align}
    h(\Value{t}{M}[s])
    & = h(\Assign{f}{M}(\Value{t_1}{M}[s], \dots, \Value{t_n}{M}[s]) \notag\\
    & = \Assign{f}{M'}(h(\Value{t_1}{M}[s]), \dots,
    h(\Value{t_n}{M}[s])) \ollabel{iso-1}\\
    & = \Assign{f}{M'}(\Value{t_1}{M'}[h \circ s], \dots,
    \Value{t_n}{M'}[h \circ s]) \ollabel{iso-2}\\
    & = \Value{t}{M'}[h\circ s] \notag
  \end{align}
  येथे \olref{iso-1} हे \olref{defn:isomorphism} मधील
  \olref{defn:iso-func} वरून, तर \olref{iso-2} हे विगमन गृहीतकावरून मिळते.
\end{enumerate}
(b) भागाचा पुरावा सरावासाठी सोडला आहे.

$!A$ हे वाक्य असेल, तर मूल्यांकने~$s$ आणि $h \circ s$ महत्त्वाची राहत
नाहीत; म्हणून $\Sat{M}{!A}$ तेव्हा आणि केवळ तेव्हाच जेव्हा
$\Sat{M'}{!A}$.
\end{proof}

\begin{prob}
\olref[mod][bas][iso]{thm:isom} मधील (b) चा पुरावा सविस्तर पूर्ण करा.
\olref[mod][bas][iso]{defn:isomorphism} मध्ये समरूपणाचे वैशिष्ट्य ठरवणाऱ्या
पाचही गुणधर्मांपैकी प्रत्येक गुणधर्म कुठे वापरला आहे, ते नोंदवण्याची
खात्री करा.
\end{prob}

\begin{defn}
रचना $\Struct{M}$ ची \emph{स्वयंरूपता} म्हणजे $\Struct{M}$ पासून
त्याच रचनेवर असलेले समरूपण होय.
\end{defn}

\begin{prob}
कोणतीही रचना $\Struct{M}$ असू द्या. जर $X$ हा $\Struct{M}$ चा परिभाष्य
उपसंच असेल आणि $h$ ही $\Struct{M}$ ची स्वयंरूपता असेल, तर $X
= \Setabs{h(x)}{x \in X}$ हे दाखवा (म्हणजे $h$ अंतर्गत $X$ अपरिवर्तित
राहतो).
\end{prob}


\end{document}
